\chapter{Supplementary Results\label{chap:append}}

\textit{This appendix collects reference material that is too dense for the main text: the full closed-loop cross-city roster tables for NAVSIM v1.1 and v2.2, a per-metric breakdown of the aggregate PDMS score into its five NAVSIM components, a data card for the NAVSIM City Splits benchmark, the full per-study hyperparameter recipes, and an index of the per-run mechanistic artifacts used in~\refsec{sec:exp-why-ssl}.}

% ============================================================================
\section{Full Cross-City Tables}
\label{app:cross-city-matrices}

The main text reads cross-city transfer through figures, compact matched tables, and the comprehensive synthesis table at the end of Chapter~5. The full run-by-run record is kept here. Tables~\ref{tab:navsim_pdms_merged_v11} and~\ref{tab:navsim_pdms_merged_v11_latf} are the canonical thesis tables under NAVSIM v1.1. They join the per-city PDMS values to the mechanistic readouts on the same row for the full checkpoint roster, split by planner so the appendix remains readable. Tables~\ref{tab:navsim_pdms_merged} and~\ref{tab:navsim_pdms_merged_latf} keep the earlier v2.2 PDMS-only scores for direct comparison with the companion paper. Together they document the full checkpoint roster and show that the headline structural findings persist across both devkit versions.

% Full NAVSIM v1.1 cross-city PDMS matrix, canonical thesis table.
\begin{sidewaystable}[p]
    \centering
    \caption[NAVSIM v1.1 cross-city PDMS matrix (canonical).]{%
    Full closed-loop cross-city PDMS matrix on NAVSIM, devkit v1.1. This is the canonical closed-loop table used for every cross-city claim in the thesis. Each row is a single trained model; \emph{Train City} marks the cities used to fit that model. The two planner blocks (\emph{TransFuser}, \emph{Latent TransFuser}) report destination-city PDMS on the four NAVSIM test splits (LV, B, P, S) plus a \emph{Total} column averaging the four.
    How to read it: rows with all four training cities checked measure mixed-city interpolation; rows with a single training city measure zero-shot geographic transfer, in which case the three non-checked columns are the out-of-distribution cells. The gap between the in-distribution column and the OOD mean under a single planner isolates the transfer penalty; comparing the same row across the two planners separates backbone effects from head effects. All values are PDMS percentages; higher is better.}
    \label{tab:navsim_pdms_merged_v11}
    \scriptsize
    \setlength{\tabcolsep}{2pt}
    \renewcommand{\arraystretch}{0.95}
    \resizebox{\textwidth}{!}{%
    \begin{tabular}{|l|cccc|ccccc|ccccc|}
    \toprule
    \multirow{2}{*}{\textbf{Backbone}}
    & \multicolumn{4}{c|}{\textbf{Train City}}
    & \multicolumn{5}{c|}{\textbf{TransFuser (PDMS \%)}}
    & \multicolumn{5}{c|}{\textbf{Latent TransFuser (PDMS \%)}} \\
    \cmidrule(lr){2-5}\cmidrule(lr){6-10}\cmidrule(lr){11-15}
    & LV & B & P & S
    & LV & B & P & S & Total
    & LV & B & P & S & Total \\
    \midrule
    ResNet-34 supervised          & \cmark & \cmark & \cmark & \cmark & 86.6 & \textbf{78.4} & \textbf{73.8} & \textbf{64.0} & \textbf{77.9} & \textbf{89.0} & \textbf{82.5} & \textbf{80.3} & \textbf{66.5} & \textbf{81.7} \\
    ResNet-34 supervised          & --     & \cmark & --     & --     & 75.3 & 73.3 & 53.6 & 41.5 & 65.0 & 67.3 & 71.7 & 54.5 & 45.0 & 62.6 \\
    ResNet-34 supervised          & \cmark & --     & --     & --     & \textbf{89.2} & 67.0 & 55.1 & 45.3 & 68.6 & 86.9 & 62.1 & 53.2 & 42.3 & 65.5 \\
    ResNet-34 supervised          & --     & --     & \cmark & --     & 68.8 & 55.3 & 52.9 & 43.0 & 57.4 & 67.7 & 57.5 & 53.0 & 46.5 & 58.3 \\
    ResNet-34 supervised          & --     & --     & --     & \cmark & 56.6 & 45.4 & 49.6 & 42.1 & 49.5 & 57.4 & 42.6 & 45.1 & 45.5 & 48.5 \\
    \midrule
    I-JEPA ViT-H/14 frozen        & \cmark & \cmark & \cmark & \cmark & 85.3 & 75.7 & 72.4 & 67.3 & 76.9 & 83.7 & 75.2 & 70.9 & \textbf{68.3} & 76.1 \\
    I-JEPA ViT-H/14 trainable     & \cmark & \cmark & \cmark & \cmark & 86.6 & \textbf{77.1} & 70.6 & 64.8 & 77.0 & \textbf{88.7} & \textbf{77.3} & 75.1 & 67.1 & \textbf{79.1} \\
    I-JEPA ViT-S/14 nuSc rect., frozen        & \cmark & \cmark & \cmark & \cmark & 87.3 & 77.0 & \textbf{74.3} & \textbf{69.2} & \textbf{78.7} & 86.6 & 73.3 & 69.7 & 59.9 & 74.9 \\
    I-JEPA ViT-S/14 nuSc rect., trainable     & \cmark & \cmark & \cmark & \cmark & 89.0 & 76.8 & 73.3 & 67.0 & 78.6 & 86.2 & 73.9 & 74.7 & 59.5 & 75.9 \\
    I-JEPA ViT-S/14 nuSc sq., frozen          & \cmark & \cmark & \cmark & \cmark & 84.8 & 75.1 & 69.7 & 69.1 & 76.3 & 83.8 & 73.2 & 68.6 & 61.5 & 74.0 \\
    I-JEPA ViT-S/14 nuSc sq., trainable       & \cmark & \cmark & \cmark & \cmark & 85.8 & 75.7 & 73.0 & 65.2 & 76.9 & 87.3 & 76.4 & \textbf{75.3} & 67.7 & 78.5 \\
    I-JEPA ViT-H/14 frozen        & --     & \cmark & --     & --     & 72.5 & 69.9 & 58.2 & 50.1 & 65.3 & 69.4 & 67.5 & 62.2 & 42.2 & 63.1 \\
    I-JEPA ViT-H/14 trainable     & --     & \cmark & --     & --     & 72.0 & 68.9 & 60.1 & 42.0 & 63.9 & 70.7 & 68.2 & 60.7 & 48.7 & 64.5 \\
    I-JEPA ViT-H/14 frozen        & \cmark & --     & --     & --     & 86.9 & 61.2 & 53.7 & 51.4 & 66.7 & 85.9 & 57.5 & 51.9 & 47.8 & 64.3 \\
    I-JEPA ViT-H/14 trainable     & \cmark & --     & --     & --     & 87.7 & 60.1 & 53.9 & 51.3 & 66.7 & 82.7 & 60.3 & 54.8 & 44.3 & 64.2 \\
    I-JEPA ViT-H/14 frozen        & --     & --     & \cmark & --     & 71.4 & 61.3 & 64.5 & 49.5 & 63.5 & 71.8 & 63.5 & 69.0 & 48.6 & 65.0 \\
    I-JEPA ViT-H/14 trainable     & --     & --     & \cmark & --     & 69.1 & 59.6 & 65.3 & 39.0 & 60.7 & 63.7 & 60.1 & 59.4 & 35.4 & 57.3 \\
    I-JEPA ViT-H/14 frozen        & --     & --     & --     & \cmark & 36.0 & 35.9 & 45.1 & 35.6 & 37.7 & 58.3 & 44.6 & 43.0 & 55.6 & 50.6 \\
    I-JEPA ViT-H/14 trainable     & --     & --     & --     & \cmark & 56.7 & 45.3 & 48.0 & 50.3 & 50.4 & 62.7 & 43.9 & 44.4 & 52.8 & 51.7 \\
    \midrule
    DINOv2 ViT-S/14 frozen                    & \cmark & \cmark & \cmark & \cmark & 87.7 & 77.0 & 71.5 & 66.0 & 77.7 & 86.4 & 72.8 & 67.2 & 59.2 & 74.1 \\
    DINOv2 ViT-S/14 trainable                 & \cmark & \cmark & \cmark & \cmark & 85.4 & 73.5 & 67.6 & 68.9 & 75.6 & 82.7 & 68.9 & 64.0 & 54.2 & 70.2 \\
    DINOv2 ViT-S/14 nuSc rect., frozen        & \cmark & \cmark & \cmark & \cmark & 86.6 & 77.2 & 74.0 & 70.5 & 78.7 & 87.1 & 74.8 & 71.6 & \textbf{71.5} & 77.7 \\
    DINOv2 ViT-S/14 nuSc rect., trainable     & \cmark & \cmark & \cmark & \cmark & 87.9 & 79.0 & 76.2 & 67.4 & \textbf{79.6} & 87.3 & \textbf{78.5} & \textbf{75.7} & 66.7 & \textbf{79.0} \\
    DINOv2 ViT-S/14 nuSc sq., frozen          & \cmark & \cmark & \cmark & \cmark & 87.2 & \textbf{79.2} & \textbf{77.8} & 65.5 & 79.4 & 86.3 & 73.1 & 72.7 & 60.8 & 75.5 \\
    DINOv2 ViT-S/14 nuSc sq., trainable       & \cmark & \cmark & \cmark & \cmark & 86.3 & 75.9 & 69.0 & \textbf{71.2} & 77.3 & 86.7 & 73.6 & 70.6 & 71.2 & 77.0 \\
    \midrule
    MAE ViT-B/16 frozen                       & \cmark & \cmark & \cmark & \cmark & 89.0 & 80.8 & \textbf{78.7} & 65.1 & 80.7 & 81.7 & 70.7 & 62.8 & 69.9 & 72.7 \\
    MAE ViT-B/16 trainable                    & \cmark & \cmark & \cmark & \cmark & 85.4 & 78.0 & 73.9 & 67.4 & 78.0 & 88.7 & \textbf{81.9} & \textbf{80.0} & 61.6 & 80.6 \\
    MAE ViT-S/14 nuSc rect., frozen           & \cmark & \cmark & \cmark & \cmark & 89.2 & 79.1 & 75.1 & \textbf{73.1} & 80.7 & 87.4 & 79.2 & 79.8 & 64.6 & 79.8 \\
    MAE ViT-S/14 nuSc rect., trainable        & \cmark & \cmark & \cmark & \cmark & 89.3 & \textbf{81.2} & 78.4 & 72.1 & \textbf{81.9} & 86.2 & 79.7 & 77.8 & 72.3 & 80.3 \\
    MAE ViT-S/14 nuSc sq., frozen             & \cmark & \cmark & \cmark & \cmark & 87.7 & 77.3 & 75.2 & 67.8 & 78.9 & 86.4 & 75.0 & 70.3 & 67.9 & 76.7 \\
    MAE ViT-S/14 nuSc sq., trainable          & \cmark & \cmark & \cmark & \cmark & 86.9 & 77.3 & 76.3 & 70.3 & 79.2 & \textbf{89.3} & 78.5 & 78.0 & \textbf{74.1} & \textbf{81.4} \\
    MAE ViT-B/16 frozen                       & --     & \cmark & --     & --     & 70.6 & 67.0 & 56.8 & 43.4 & 62.4 & 71.8 & 73.4 & 60.0 & 48.7 & 66.3 \\
    MAE ViT-B/16 trainable                    & --     & \cmark & --     & --     & 70.6 & 67.3 & 58.0 & 41.9 & 62.5 & 69.4 & 73.1 & 61.3 & 53.7 & 66.5 \\
    MAE ViT-B/16 frozen                       & \cmark & --     & --     & --     & 88.3 & 66.2 & 55.2 & 52.9 & 69.3 & 86.6 & 66.0 & 56.1 & 52.2 & 68.8 \\
    MAE ViT-B/16 trainable                    & \cmark & --     & --     & --     & 87.4 & 59.8 & 53.2 & 51.6 & 66.4 & 87.8 & 63.3 & 56.0 & 51.0 & 68.1 \\
    MAE ViT-B/16 frozen                       & --     & --     & \cmark & --     & 69.3 & 59.9 & 61.1 & 48.9 & 61.5 & 71.2 & 54.8 & 50.4 & 52.6 & 59.1 \\
    MAE ViT-B/16 trainable                    & --     & --     & \cmark & --     & 73.0 & 65.6 & 68.8 & 40.8 & 64.8 & 64.4 & 63.9 & 61.0 & 35.7 & 59.1 \\
    MAE ViT-B/16 frozen                       & --     & --     & --     & \cmark & 58.0 & 45.3 & 51.2 & 56.7 & 52.5 & 54.6 & 44.0 & 46.6 & 49.7 & 49.0 \\
    MAE ViT-B/16 trainable                    & --     & --     & --     & \cmark & 51.0 & 46.0 & 52.2 & 36.0 & 47.3 & 55.6 & 43.8 & 48.0 & 42.5 & 48.4 \\
    \bottomrule
    \end{tabular}}
\end{sidewaystable}

\begin{table}[ht]
\centering
\caption[Boston-only v1.1 Latent TransFuser cross-city summary.]{%
Compact summary of the canonical Boston-trained Latent TransFuser rows from the v1.1 matrix in~\reftab{tab:navsim_pdms_merged_v11}. These four rows are used as the matched supervised-vs-SSL comparison throughout the thesis: each model is trained on Boston only and evaluated on every NAVSIM city. LV/B/P/S are destination-city PDMS percentages; the in-distribution column (bold) is Boston; \emph{OOD mean} averages the three remaining cities and \emph{Drop} is the in-distribution-minus-OOD-mean gap in percentage points. A smaller drop indicates a more geographically robust representation under a fixed planner. The MAE row has the best absolute OOD mean; the DINOv2 row has the smallest drop; the ResNet-34 row has the largest drop and is the supervised reference.}
\label{tab:navsim_pdms_boston_v11}
\vspace{0.35em}
\small
\renewcommand{\arraystretch}{1.15}
\setlength{\tabcolsep}{5pt}
\begin{tabular}{@{}>{\raggedright\arraybackslash}p{3.6cm} c c c c c c@{}}
\toprule
\textbf{Backbone} & \textbf{LV (\%)} & \textbf{B (\%)} & \textbf{P (\%)} & \textbf{S (\%)} & \textbf{OOD mean (\%)} & \textbf{Drop (pp)} \\
\midrule
\rowcolor{gray!12}
ResNet-34 supervised & 67.3 & \textbf{71.7} & 54.5 & 45.0 & 55.6 & \textbf{16.1} \\
MAE ViT-B/16 frozen  & 71.8 & \textbf{73.4} & 60.0 & 48.7 & \textbf{60.2} & 13.2 \\
I-JEPA ViT-H/14 frozen & 69.4 & \textbf{67.5} & 62.2 & 42.2 & 57.9 & 9.6 \\
DINOv2 ViT-S/14 frozen & 70.7 & \textbf{63.6} & 55.8 & 45.5 & 57.3 & \textbf{6.3} \\
\bottomrule
\end{tabular}
\end{table}

% Full NAVSIM v2.2 cross-city PDMS matrix (legacy devkit).
% Retained as provenance for the companion paper's numbers; the canonical
% thesis matrix is the v1.1 table (tab_navsim_pdms_merged_v11.tex).
\begin{sidewaystable}[p]
    \centering
    \caption[NAVSIM v2.2 cross-city PDMS matrix (legacy).]{%
    Full closed-loop cross-city PDMS matrix on NAVSIM, devkit v2.2 (legacy).
    Each row is a single trained model. The left block (\emph{Train City}) marks which of the four cities were used for training; \cmark indicates inclusion. Total combines performance across all four evaluation cities for each planner. All numbers are PDMS (\%); higher is better.
    How to read it: rows with four training cities are the all-city baselines, and their \emph{Total} column is mixed-city interpolation performance. Rows with a single training city are zero-shot transfer probes; non-checked columns under each planner are the out-of-distribution cells, and the gap between the in-distribution column and the three OOD columns quantifies the geographic transfer penalty. This matrix is retained for direct comparison with the companion cross-city benchmark paper~\cite{naeinian-cross-city-2026}; the thesis claims use the v1.1 re-evaluation in~\reftab{tab:navsim_pdms_merged_v11}.}
    \label{tab:navsim_pdms_merged}
    \scriptsize
    \setlength{\tabcolsep}{2pt}
    \renewcommand{\arraystretch}{0.95}
    \resizebox{\textwidth}{!}{%
    \begin{tabular}{|l|cccc|ccccc|ccccc|}
    \toprule
    \multirow{2}{*}{\textbf{Backbone}}
    & \multicolumn{4}{c|}{\textbf{Train City}}
    & \multicolumn{5}{c|}{\textbf{TransFuser (PDMS \%)}}
    & \multicolumn{5}{c|}{\textbf{Latent TransFuser (PDMS \%)}} \\
    \cmidrule(lr){2-5}\cmidrule(lr){6-10}\cmidrule(lr){11-15}
    & LV & B & P & S
    & LV & B & P & S & Total
    & LV & B & P & S & Total \\
    \midrule
    ResNet-34 supervised          & \cmark & \cmark & \cmark & \cmark & 88.4 & \textbf{81.4} & \textbf{75.7} & \textbf{59.9} & \textbf{79.2} & 87.0 & \textbf{79.3} & \textbf{74.4} & \textbf{67.2} & \textbf{79.0} \\
    ResNet-34 supervised          & --     & \cmark & --     & --     & 75.3 & 73.3 & 53.6 & 41.5 & 65.0 & 67.9 & 72.2 & 50.0 & 44.0 & 61.8 \\
    ResNet-34 supervised          & \cmark & --     & --     & --     & \textbf{89.2} & 67.0 & 55.1 & 45.3 & 68.6 & \textbf{87.2} & 63.6 & 50.9 & 41.9 & 65.5 \\
    ResNet-34 supervised          & --     & --     & \cmark & --     & 68.8 & 55.3 & 52.9 & 43.0 & 57.4 & 68.2 & 58.0 & 51.7 & 44.3 & 58.0 \\
    ResNet-34 supervised          & --     & --     & --     & \cmark & 57.9 & 45.5 & 43.4 & 41.0 & 48.5 & 59.3 & 43.9 & 41.1 & 44.6 & 48.6 \\
    \midrule
    I-JEPA ViT-H/14 frozen        & \cmark & \cmark & \cmark & \cmark & 86.5 & 77.1 & 71.2 & 68.1 & 77.6 & 85.0 & 76.5 & 69.6 & 68.7 & 76.7 \\
    I-JEPA ViT-H/14 100\% trainable & \cmark & \cmark & \cmark & \cmark & 87.4 & \textbf{77.7} & 68.8 & 64.2 & 77.0 & \textbf{89.3} & \textbf{77.8} & \textbf{73.0} & 67.6 & \textbf{79.1} \\
    I-JEPA ViT-S/14 nuSc rect., frozen        & \cmark & \cmark & \cmark & \cmark & 88.4 & 77.5 & \textbf{71.5} & 68.4 & \textbf{78.5} & 87.6 & 76.5 & 71.1 & 57.0 & 76.1 \\
    I-JEPA ViT-S/14 nuSc rect., trainable     & \cmark & \cmark & \cmark & \cmark & 88.9 & \textbf{77.7} & 70.8 & 67.0 & 78.4 & 86.6 & 74.8 & 67.2 & 65.3 & 75.7 \\
    I-JEPA ViT-S/14 nuSc sq., frozen          & \cmark & \cmark & \cmark & \cmark & 85.5 & 75.9 & 68.2 & \textbf{69.4} & 76.6 & 84.5 & 73.3 & 68.0 & 65.1 & 74.7 \\
    I-JEPA ViT-S/14 nuSc sq., trainable       & \cmark & \cmark & \cmark & \cmark & 86.7 & 76.2 & 70.9 & 65.7 & 77.0 & 82.9 & 69.8 & 59.1 & \textbf{70.2} & 72.1 \\
    I-JEPA ViT-H/14 frozen        & --     & \cmark & --     & --     & 74.1 & 71.2 & 54.2 & 48.8 & 65.2 & 70.2 & 66.4 & 54.7 & 41.2 & 61.3 \\
    I-JEPA ViT-H/14 trainable     & --     & \cmark & --     & --     & 72.5 & 68.0 & 52.4 & 40.9 & 62.1 & 71.7 & 68.5 & 55.0 & 47.8 & 63.6 \\
    I-JEPA ViT-H/14 frozen        & \cmark & --     & --     & --     & 88.4 & 63.6 & 52.5 & 50.5 & 67.6 & 87.2 & 59.2 & 50.5 & 46.4 & 64.8 \\
    I-JEPA ViT-H/14 trainable     & \cmark & --     & --     & --     & 88.3 & 62.2 & 52.9 & 50.1 & 67.1 & 83.1 & 62.0 & 52.5 & 44.7 & 64.4 \\
    I-JEPA ViT-H/14 frozen        & --     & --     & \cmark & --     & 72.1 & 62.2 & 62.5 & 47.4 & 63.3 & 72.0 & 62.9 & 63.9 & 47.5 & 63.7 \\
    I-JEPA ViT-H/14 trainable     & --     & --     & \cmark & --     & 70.1 & 60.4 & 60.1 & 37.4 & 60.0 & 64.0 & 60.1 & 55.8 & 33.8 & 56.4 \\
    I-JEPA ViT-H/14 frozen        & --     & --     & --     & \cmark & 35.1 & 30.5 & 34.1 & 33.4 & 33.2 & 59.6 & 45.1 & 38.8 & 53.9 & 50.1 \\
    I-JEPA ViT-H/14 trainable     & --     & --     & --     & \cmark & 58.1 & 44.9 & 43.8 & 48.8 & 49.7 & 64.0 & 44.9 & 41.3 & 50.8 & 51.5 \\
    \midrule
    DINOv2 ViT-S/14 frozen                    & \cmark & \cmark & \cmark & \cmark & 88.7 & 78.4 & 70.2 & 66.1 & 78.3 & 87.6 & 76.5 & 69.8 & 62.2 & 76.6 \\
    DINOv2 ViT-S/14 trainable                 & \cmark & \cmark & \cmark & \cmark & 86.0 & 74.4 & 66.0 & 68.9 & 75.7 & 86.4 & 72.1 & 63.2 & 56.0 & 72.5 \\
    DINOv2 ViT-S/14 nuSc rect., frozen        & \cmark & \cmark & \cmark & \cmark & 87.2 & 78.1 & 72.1 & 71.0 & 78.8 & 87.9 & 75.5 & 70.2 & 66.6 & 77.2 \\
    DINOv2 ViT-S/14 nuSc rect., trainable     & \cmark & \cmark & \cmark & \cmark & 88.5 & 79.7 & 74.1 & 67.9 & \textbf{79.7} & \textbf{88.3} & \textbf{79.8} & \textbf{73.1} & \textbf{71.3} & \textbf{79.9} \\
    DINOv2 ViT-S/14 nuSc sq., frozen          & \cmark & \cmark & \cmark & \cmark & 88.0 & \textbf{80.4} & \textbf{75.9} & 65.8 & \textbf{79.7} & 86.8 & 75.0 & 69.5 & 61.2 & 75.7 \\
    DINOv2 ViT-S/14 nuSc sq., trainable       & \cmark & \cmark & \cmark & \cmark & 87.7 & 77.1 & 68.4 & \textbf{71.5} & 78.0 & 83.7 & 73.0 & 64.2 & 66.7 & 73.8 \\
    \midrule
    MAE ViT-B/16 frozen                       & \cmark & \cmark & \cmark & \cmark & 89.5 & 80.9 & 75.1 & 65.5 & 80.2 & 87.6 & 76.9 & 70.5 & 69.5 & 78.0 \\
    MAE ViT-B/16 trainable                    & \cmark & \cmark & \cmark & \cmark & 86.3 & 78.6 & 69.3 & 67.8 & 77.6 & 86.3 & 76.5 & 68.5 & \textbf{74.9} & 77.9 \\
    MAE ViT-S/14 nuSc rect., frozen           & \cmark & \cmark & \cmark & \cmark & \textbf{90.3} & 80.2 & 74.2 & \textbf{73.7} & 81.4 & 85.3 & 77.3 & 72.9 & 69.9 & 77.9 \\
    MAE ViT-S/14 nuSc rect., trainable        & \cmark & \cmark & \cmark & \cmark & 89.9 & \textbf{82.4} & \textbf{77.4} & 72.6 & \textbf{82.4} & 89.2 & \textbf{79.9} & 70.2 & 68.9 & \textbf{79.3} \\
    MAE ViT-S/14 nuSc sq., frozen             & \cmark & \cmark & \cmark & \cmark & 88.9 & 78.6 & 73.7 & 68.3 & 79.4 & 84.5 & 75.2 & 71.5 & 67.3 & 76.3 \\
    MAE ViT-S/14 nuSc sq., trainable          & \cmark & \cmark & \cmark & \cmark & 87.4 & 77.8 & 74.2 & 70.6 & 79.2 & 85.8 & 77.7 & \textbf{74.4} & 64.1 & 77.6 \\
    \bottomrule
    \end{tabular}}%

    \vspace{0.6em}
    \footnotesize
    \emph{Single-city transfer rows for DINOv2 and MAE (Boston, Pittsburgh, Las~Vegas, Singapore) follow the same convention and are reported in the companion paper's supplementary material~\cite{naeinian-cross-city-2026}; the reduced set above is retained as the minimum provenance needed to justify the v1.1 re-evaluation used for all thesis claims.}
\end{sidewaystable}


% ============================================================================
\section{Per-Metric PDMS Decomposition}
\label{app:per-metric}

NAVSIM's aggregate PDMS is a weighted combination of five sub-scores: no-at-fault collisions (NC), drivable-area compliance (DAC), ego progress (EP), time-to-collision (TTC), and comfort. The cross-city matrices in the main text report only the aggregate, which is the primary robustness indicator used by every thesis claim. A component-level decomposition is still useful, because it isolates whether cross-city degradation is dominated by geometry-aware sub-scores (DAC, EP) or by agent-interaction sub-scores (NC, TTC). For the completed DiffusionDrive rows, the answer is clear: Singapore transfer failures concentrate in DAC and EP, while NC and TTC remain comparatively stable.

On structural grounds, the expected pattern under city shift is:

\begin{itemize}
    \item \textbf{DAC and EP} should degrade most, since both depend on lane geometry, lane width, and turn radii, all of which differ systematically between the four cities.
    \item \textbf{NC and TTC} should be more robust, since they depend primarily on agent-to-agent spacing, which is more invariant across cities than lane geometry.
    \item \textbf{Comfort} should be weakly city-dependent and dominated by speed-profile smoothness.
\end{itemize}

\begin{table}[ht]
\centering
\caption[Singapore component breakdown for completed cross-city DiffusionDrive rows.]{Singapore destination component scores for the three completed cross-city DiffusionDrive rows. Each row is one trained model evaluated on Singapore. Values are percentages for the five PDMS sub-scores used in the main-text aggregate. The matched Pittsburgh pair is the load-bearing comparison: frozen I-JEPA improves drivable-area compliance on Singapore while leaving the collision-adjacent metrics essentially unchanged.}
\label{tab:dd-cc-singapore-components}
\small
\renewcommand{\arraystretch}{1.15}
\setlength{\tabcolsep}{5pt}
\begin{tabular}{@{}>{\raggedright\arraybackslash}p{3.4cm} c c c c c@{}}
\toprule
\textbf{Train $\rightarrow$ test} & \textbf{NC (\%)} & \textbf{DAC (\%)} & \textbf{EP (\%)} & \textbf{TTC (\%)} & \textbf{Comfort (\%)} \\
\midrule
Rn34-bos $\rightarrow$ Singapore & 95.5 & 62.8 & 79.0 & 94.7 & 45.4 \\
Rn34-pit $\rightarrow$ Singapore & 96.3 & 61.7 & 73.7 & 95.4 & 47.8 \\
IJSr-pit $\rightarrow$ Singapore & 96.7 & 64.1 & 72.2 & 96.0 & 50.4 \\
\bottomrule
\end{tabular}
\end{table}

The matched Pittsburgh pair is the clearest read. Frozen I-JEPA improves DAC by $2.4$ points on Singapore relative to the supervised ResNet-34 row, while NC and TTC remain within $0.6$ points. The representation gain under cross-city DiffusionDrive therefore shows up in geometry-sensitive compliance rather than in a wholesale change to reactive safety metrics.

% ============================================================================
\section{NAVSIM City Splits Data Card}
\label{app:city-data-card}

This data card documents the scene-count imbalance inherent in the nuPlan-derived NAVSIM city splits. The imbalance is directly relevant to the Singapore result in the cross-city benchmark and to the interpretation of the mechanistic effective-rank collapses that cluster on Singapore-trained runs.

% City-split data card. Counts taken from navsim-ssl-city-generalization/city_splits/splits/*.json
% (derived from nuPlan map_location tags; see navsim-ssl-city-generalization/city_splits/create_splits.py).
\begin{table}[ht]
    \centering
    \caption[NAVSIM city-split data card.]{%
    NAVSIM city-split data card. Logs are grouped by nuPlan city tag. The \emph{trainval} column is used for training and validation; the \emph{test} column is the held-out \texttt{navtest} subset used for cross-city evaluation. The \texttt{map\_location} column records the underlying nuPlan map identifier.%
    }
    \label{tab:city-splits}
    \small
    \begin{tabular}{lccl}
    \toprule
    \textbf{City} & \textbf{Trainval logs} & \textbf{Test logs} & \textbf{nuPlan \texttt{map\_location}} \\
    \midrule
    Las Vegas (LV)   & 818 & 50 & \texttt{us-nv-las-vegas-strip} \\
    Boston (B)       & 255 & 62 & \texttt{us-ma-boston} \\
    Pittsburgh (P)   & 140 & 22 & \texttt{us-pa-pittsburgh-hazelwood} \\
    Singapore (S)    &  97 & 13 & \texttt{sg-one-north} \\
    \midrule
    \textbf{Total}   & \textbf{1310} & \textbf{147} & -- \\
    \bottomrule
    \end{tabular}
\end{table}


The imbalance is severe: Las Vegas contributes $818$ training logs against Singapore's $97$. This is a consequence of the original nuPlan capture plan rather than a decision on our part; the NAVSIM benchmark inherits the imbalance directly. For every cross-city benchmark single-city result where the Singapore-trained row underperforms, the reader should ask whether the effect reflects a representation difference or a data-scale difference. The mechanistic feature-rank analysis (\refsec{sec:exp-implicit-reg}) finds that the three most extreme feature-rank collapses involve either trainable ViTs on all-city data (not a small-data regime) or frozen ViTs on single-city non-Singapore runs, so the small-data hypothesis is not a clean explanation. Still, the imbalance is a confounder the reader should hold in mind; disentangling it fully requires a data-efficiency study that is out of scope for this thesis.

% ============================================================================
\section{Per-Study Hyperparameter Tables}
\label{app:hparams}

The recipes below are intended to be sufficient for reproduction. The NAVSIM evaluation devkit version is v2.2 for the numbers that match the companion cross-city benchmark paper and v1.1 for the canonical thesis matrix in~\reftab{tab:navsim_pdms_merged_v11}.

% Per-study training recipes used throughout the thesis.
\begin{table}[ht]
    \centering
    \caption[Training recipe: LAW (cross-city benchmark).]{Training recipe for the LAW agent in the cross-city benchmark's open-loop nuScenes evaluation. LAW rows vary only by training-city filter and random seed. The remaining settings follow the LAW reference implementation~\cite{li-law-2024}.}
    \label{tab:hparam-phase1-law}
    \small
    \begin{tabular}{ll}
    \toprule
    Optimizer             & AdamW \\
    Learning rate         & $5\times10^{-5}$ \\
    Weight decay          & $0.01$ \\
    Scheduler             & Cosine, no warm-up \\
    Batch size            & 2 (effective) \\
    Epochs                & 12 \\
    Gradient clip         & 35 \\
    GPUs                  & 1 $\times$ L40S \\
    Precision             & 16-mixed \\
    \bottomrule
    \end{tabular}
\end{table}

\begin{table}[ht]
    \centering
    \caption[Training recipe: TransFuser / Latent TransFuser.]{Training recipe for the TransFuser and Latent TransFuser v1.1 cross-city matrix in~\reftab{tab:navsim_pdms_merged_v11}. The recipe follows the original TransFuser setup: plain Adam, no weight decay, no learning-rate schedule, and gradient clipping at $1.0$.}
    \label{tab:hparam-phase1-tf}
    \small
    \begin{tabular}{ll}
    \toprule
    Optimizer             & plain Adam (\textbf{not} AdamW) \\
    Learning rate         & $1\times10^{-4}$ \\
    Weight decay          & \textbf{0} (no decay) \\
    Scheduler             & \textbf{None} (constant lr) \\
    Batch size            & 32/GPU, effective 128 \\
    Epochs                & 20 (paper baseline) or 100 (full) \\
    Gradient clip         & 1.0 \\
    LiDAR backbone        & ResNet-34 (\textbf{not} ResNet-18) \\
    GPUs                  & 4 $\times$ H200 \\
    Precision             & 16-mixed \\
    \bottomrule
    \end{tabular}
\end{table}

\begin{table}[ht]
    \centering
    \caption[Training recipe: DiffusionDrive + SSL (generative-planning study).]{Training recipe for the DiffusionDrive$+$SSL rows of the generative-planning study. The image encoder is frozen by default, the trajectory head uses a $256$-anchor K-means bank, and inference uses two denoising steps. Other settings follow the base DiffusionDrive recipe~\cite{jia-diffusiondrive-2025}.}
    \label{tab:hparam-phase2-dd}
    \small
    \begin{tabular}{ll}
    \toprule
    Optimizer             & Adam \\
    Learning rate         & $1\times10^{-4}$ \\
    Weight decay          & $0$ \\
    Scheduler             & Constant \\
    Batch size            & 32/GPU, effective 128 \\
    Epochs                & 20 \\
    K-means anchor count  & 256 \\
    Denoising steps       & 2 (truncated) \\
    GPUs                  & 2 $\times$ H200 \\
    Precision             & 16-mixed \\
    \bottomrule
    \end{tabular}
\end{table}

\begin{table}[ht]
    \centering
    \caption[Pre-training recipe: I-JEPA ViT-S/14 on nuScenes.]{Domain-specific self-supervised pre-training recipe for I-JEPA ViT-S/14 on nuScenes frames. The rectangular variant uses $224 \times 392$ input, and the square variant uses $224 \times 224$ input. The DINOv2 and MAE domain-specific variants use the same substrate with the objective swapped.}
    \label{tab:hparam-phase3-ssl-pretrain}
    \small
    \begin{tabular}{ll}
    \toprule
    Architecture          & ViT-S/14, 22M parameters \\
    Objective             & I-JEPA target-vs-context prediction \\
    Optimizer             & AdamW \\
    Learning rate         & Cosine schedule, peak $5\times10^{-4}$ \\
    Weight decay          & $0.04$ \\
    Epochs                & 100 \\
    Batch size            & 64/GPU, effective 128 \\
    Input resolution      & $224\times224$ (sq.) or $224\times392$ (rect.) \\
    Dataset               & nuScenes front-camera frames (340k images) \\
    GPUs                  & 2 $\times$ H200 \\
    Precision             & 16-mixed \\
    \bottomrule
    \end{tabular}
\end{table}


A subtle but load-bearing detail: the TransFuser recipe in~\reftab{tab:hparam-phase1-tf} uses plain Adam with no weight decay and no learning-rate schedule. Substituting AdamW, or adding a cosine schedule, silently regresses the in-distribution PDMS by more than thirty points with no accompanying error signal. This was the root cause of an early training regression during the cross-city benchmark build-out, and the recipe above is the one that reproduces the published TransFuser mean.

% ============================================================================
\section{Mechanistic Artifacts Index}
\label{app:mechanistic-index}

The mechanistic analysis in~\refsec{sec:exp-why-ssl} produces one CKA heatmap and one singular-value spectrum per analysed checkpoint. The table below indexes every analysed run by its label, pairs it with the family-level probe, CKA, and effective-rank numbers that the analysis aggregates over, and lets a reader tracing an individual row of~\reftab{tab:navsim_pdms_merged_v11} or~\reftab{tab:navsim_pdms_merged_v11_latf} locate the corresponding mechanistic artifacts in constant time.

% Auto-generated from fig/phase4/all_metrics.json. 188 analyzed runs.
% Each row points to per-run CKA heatmap + SVD spectrum already present
% in fig/phase4/heatmaps/ and fig/phase4/spectra/.
\begin{longtable}{lllllrrr}
\caption{Per-run mechanistic artifacts index. ``train'' uses the short city codes (LV, B, P, S) or \textit{all} for LV+B+P+S. Each run has a corresponding CKA heatmap at \texttt{fig/phase4/heatmaps/cka\_<run\_id>.pdf} and singular-value spectrum at \texttt{fig/phase4/spectra/spectrum\_<run\_id>.pdf}. Probe, CKA, and effective-rank columns summarise the per-run aggregates reported in Sec.~\ref{sec:exp-why-ssl}.} \label{tab:mechanistic-index} \\
\toprule
run\_id & backbone & variant & mode & train & probe & CKA & rank \\
\midrule
\endfirsthead
\multicolumn{8}{c}{\textit{Table~\ref{tab:mechanistic-index} continued.}} \\
\toprule
run\_id & backbone & variant & mode & train & probe & CKA & rank \\
\midrule
\endhead
\midrule
\multicolumn{8}{r}{\textit{continued on next page}} \\
\endfoot
\bottomrule
\endlastfoot
A3-b & resnet34 & supervised & TF & all & 0.99 & 0.70 & 72.0 \\
\midrule
A6 & ijepa & ViT-H/14 frozen & TF & all & 0.94 & 0.69 & 36.6 \\
A6-b & ijepa & ViT-H/14 trainable & TF & all & 0.38 & 0.70 & 2.9 \\
A6-d & ijepa & ViT-S/14 rect frozen & TF & all & 0.79 & 0.91 & 17.4 \\
A6-e & ijepa & ViT-S/14 rect trainable & TF & all & 0.78 & 0.90 & 18.5 \\
A6-f & ijepa & ViT-S/14 sq frozen & TF & all & 0.94 & 0.67 & 107.9 \\
A6-g & ijepa & ViT-S/14 sq trainable & TF & all & 0.95 & 0.57 & 90.7 \\
\midrule
A7 & dinov2 & ViT-S/14 frozen & TF & all & 0.66 & 0.95 & 4.1 \\
A7-b & dinov2 & ViT-S/14 trainable & TF & all & 0.39 & 1.00 & 1.2 \\
A7-d & dinov2 & ViT-S/14 rect frozen & TF & all & 0.85 & 0.87 & 27.3 \\
A7-e & dinov2 & ViT-S/14 rect trainable & TF & all & 0.89 & 0.88 & 25.3 \\
A7-f & dinov2 & ViT-S/14 sq frozen & TF & all & 0.87 & 0.73 & 89.3 \\
A7-g & dinov2 & ViT-S/14 sq trainable & TF & all & 0.88 & 0.71 & 95.5 \\
\midrule
A8 & mae & ViT-B/16 frozen & TF & all & 0.96 & 0.73 & 45.5 \\
A8-b & mae & ViT-B/16 trainable & TF & all & 0.87 & 0.70 & 10.7 \\
A8-d & mae & ViT-S/14 rect frozen & TF & all & 0.96 & 0.82 & 25.8 \\
A8-e & mae & ViT-S/14 rect trainable & TF & all & 0.96 & 0.56 & 21.6 \\
A8-f & mae & ViT-S/14 sq frozen & TF & all & 0.93 & 0.60 & 106.8 \\
A8-g & mae & ViT-S/14 sq trainable & TF & all & 0.98 & 0.56 & 90.5 \\
\midrule
A9 & resnet34 & supervised & LatTF & all & 1.00 & 0.62 & 69.5 \\
\midrule
A9-b & ijepa & ViT-H/14 frozen & LatTF & all & 0.94 & 0.72 & 39.5 \\
A9-c & ijepa & ViT-H/14 trainable & LatTF & all & 0.92 & 0.73 & 26.7 \\
\midrule
A9-e & dinov2 & ViT-S/14 frozen & LatTF & all & 0.83 & 0.89 & 28.2 \\
A9-f & dinov2 & ViT-S/14 trainable & LatTF & all & 0.83 & 0.79 & 16.4 \\
\midrule
A9-h & mae & ViT-B/16 frozen & LatTF & all & 0.96 & 0.76 & 41.3 \\
A9-i & mae & ViT-B/16 trainable & LatTF & all & 0.98 & 0.67 & 26.3 \\
\midrule
A9-j & ijepa & ViT-S/14 rect frozen & LatTF & all & 0.83 & 0.88 & 30.0 \\
A9-k & ijepa & ViT-S/14 rect trainable & LatTF & all & 0.90 & 0.81 & 37.9 \\
\midrule
A9-l & dinov2 & ViT-S/14 rect frozen & LatTF & all & 0.89 & 0.81 & 45.8 \\
A9-m & dinov2 & ViT-S/14 rect trainable & LatTF & all & 0.93 & 0.81 & 34.4 \\
\midrule
A9-n & mae & ViT-S/14 rect frozen & LatTF & all & 0.96 & 0.70 & 44.4 \\
A9-o & mae & ViT-S/14 rect trainable & LatTF & all & 0.97 & 0.74 & 23.9 \\
\midrule
A9-p & ijepa & ViT-S/14 sq frozen & LatTF & all & 0.92 & 0.60 & 118.6 \\
A9-q & ijepa & ViT-S/14 sq trainable & LatTF & all & 0.95 & 0.60 & 85.3 \\
\midrule
A9-r & dinov2 & ViT-S/14 sq frozen & LatTF & all & 0.88 & 0.73 & 92.5 \\
A9-s & dinov2 & ViT-S/14 sq trainable & LatTF & all & 0.89 & 0.72 & 88.6 \\
\midrule
A9-t & mae & ViT-S/14 sq frozen & LatTF & all & 0.95 & 0.63 & 104.9 \\
A9-u & mae & ViT-S/14 sq trainable & LatTF & all & 0.97 & 0.54 & 93.5 \\
\midrule
F1-B1 & resnet34 & supervised & LatTF & B & 0.97 & 0.83 & 55.0 \\
F1-B2 & resnet34 & supervised & LatTF & LV & 0.99 & 0.72 & 65.9 \\
F1-B3 & resnet34 & supervised & LatTF & P & 0.96 & 0.76 & 66.2 \\
F1-B4 & resnet34 & supervised & LatTF & S & 0.95 & 0.80 & 59.8 \\
\midrule
F1-D1 & dinov2 & ViT-S/14 frozen & LatTF & B & 0.81 & 0.93 & 8.8 \\
F1-D1-b & dinov2 & ViT-S/14 trainable & LatTF & B & 0.56 & 0.86 & 105.4 \\
F1-D1-c & dinov2 & ViT-S/14 rect frozen & LatTF & B & 0.82 & 0.90 & 30.6 \\
F1-D1-d & dinov2 & ViT-S/14 rect trainable & LatTF & B & 0.81 & 0.92 & 33.5 \\
F1-D1-e & dinov2 & ViT-S/14 sq frozen & LatTF & B & 0.76 & 0.70 & 125.8 \\
F1-D1-f & dinov2 & ViT-S/14 sq trainable & LatTF & B & 0.76 & 0.72 & 115.2 \\
F1-D2 & dinov2 & ViT-S/14 frozen & LatTF & LV & 0.84 & 0.91 & 31.5 \\
F1-D2-b & dinov2 & ViT-S/14 trainable & LatTF & LV & 0.74 & 0.87 & 12.1 \\
F1-D2-c & dinov2 & ViT-S/14 rect frozen & LatTF & LV & 0.84 & 0.82 & 43.6 \\
F1-D2-d & dinov2 & ViT-S/14 rect trainable & LatTF & LV & 0.87 & 0.78 & 46.1 \\
F1-D2-e & dinov2 & ViT-S/14 sq frozen & LatTF & LV & 0.83 & 0.71 & 105.9 \\
F1-D2-f & dinov2 & ViT-S/14 sq trainable & LatTF & LV & 0.83 & 0.75 & 97.5 \\
F1-D3 & dinov2 & ViT-S/14 frozen & LatTF & P & 0.73 & 0.90 & 5.2 \\
F1-D3-b & dinov2 & ViT-S/14 trainable & LatTF & P & 0.59 & 0.85 & 101.6 \\
F1-D3-c & dinov2 & ViT-S/14 rect frozen & LatTF & P & 0.77 & 0.88 & 36.3 \\
F1-D3-d & dinov2 & ViT-S/14 rect trainable & LatTF & P & 0.79 & 0.92 & 20.9 \\
F1-D3-e & dinov2 & ViT-S/14 sq frozen & LatTF & P & 0.74 & 0.73 & 110.1 \\
F1-D3-f & dinov2 & ViT-S/14 sq trainable & LatTF & P & 0.74 & 0.70 & 128.4 \\
F1-D4 & dinov2 & ViT-S/14 frozen & LatTF & S & 0.61 & 0.82 & 119.0 \\
F1-D4-b & dinov2 & ViT-S/14 trainable & LatTF & S & 0.56 & 0.85 & 105.6 \\
F1-D4-c & dinov2 & ViT-S/14 rect frozen & LatTF & S & 0.84 & 0.91 & 16.2 \\
F1-D4-d & dinov2 & ViT-S/14 rect trainable & LatTF & S & 0.80 & 0.95 & 9.5 \\
F1-D4-e & dinov2 & ViT-S/14 sq frozen & LatTF & S & 0.76 & 0.80 & 79.5 \\
F1-D4-f & dinov2 & ViT-S/14 sq trainable & LatTF & S & 0.77 & 0.81 & 64.1 \\
\midrule
F1-I1 & ijepa & ViT-H/14 frozen & LatTF & B & 0.83 & 0.79 & 40.2 \\
F1-I1-b & ijepa & ViT-H/14 trainable & LatTF & B & 0.77 & 0.94 & 5.5 \\
F1-I1-c & ijepa & ViT-S/14 rect frozen & LatTF & B & 0.79 & 0.90 & 24.6 \\
F1-I1-d & ijepa & ViT-S/14 rect trainable & LatTF & B & 0.81 & 0.91 & 23.3 \\
F1-I1-e & ijepa & ViT-S/14 sq frozen & LatTF & B & 0.81 & 0.66 & 163.8 \\
F1-I1-f & ijepa & ViT-S/14 sq trainable & LatTF & B & 0.82 & 0.74 & 128.1 \\
F1-I2 & ijepa & ViT-H/14 frozen & LatTF & LV & 0.87 & 0.72 & 50.2 \\
F1-I2-b & ijepa & ViT-H/14 trainable & LatTF & LV & 0.84 & 0.84 & 7.2 \\
F1-I2-c & ijepa & ViT-S/14 rect frozen & LatTF & LV & 0.80 & 0.81 & 44.0 \\
F1-I2-d & ijepa & ViT-S/14 rect trainable & LatTF & LV & 0.86 & 0.79 & 40.3 \\
F1-I2-e & ijepa & ViT-S/14 sq frozen & LatTF & LV & 0.87 & 0.68 & 119.2 \\
F1-I2-f & ijepa & ViT-S/14 sq trainable & LatTF & LV & 0.87 & 0.72 & 95.7 \\
F1-I3 & ijepa & ViT-H/14 frozen & LatTF & P & 0.81 & 0.80 & 29.1 \\
F1-I3-b & ijepa & ViT-H/14 trainable & LatTF & P & 0.76 & 0.99 & 5.5 \\
F1-I3-c & ijepa & ViT-S/14 rect frozen & LatTF & P & 0.76 & 0.83 & 27.8 \\
F1-I3-d & ijepa & ViT-S/14 rect trainable & LatTF & P & 0.77 & 0.87 & 26.4 \\
F1-I3-e & ijepa & ViT-S/14 sq frozen & LatTF & P & 0.81 & 0.68 & 158.0 \\
F1-I3-f & ijepa & ViT-S/14 sq trainable & LatTF & P & 0.83 & 0.74 & 105.4 \\
F1-I4 & ijepa & ViT-H/14 frozen & LatTF & S & 0.81 & 0.86 & 16.7 \\
F1-I4-b & ijepa & ViT-H/14 trainable & LatTF & S & 0.73 & 0.98 & 10.1 \\
F1-I4-c & ijepa & ViT-S/14 rect frozen & LatTF & S & 0.76 & 0.91 & 13.7 \\
F1-I4-d & ijepa & ViT-S/14 rect trainable & LatTF & S & 0.80 & 0.92 & 15.9 \\
F1-I4-e & ijepa & ViT-S/14 sq frozen & LatTF & S & 0.81 & 0.71 & 159.3 \\
F1-I4-f & ijepa & ViT-S/14 sq trainable & LatTF & S & 0.82 & 0.81 & 84.2 \\
\midrule
F1-M1 & mae & ViT-B/16 frozen & LatTF & B & 0.89 & 0.86 & 53.5 \\
F1-M1-b & mae & ViT-B/16 trainable & LatTF & B & 0.89 & 0.83 & 21.3 \\
F1-M1-c & mae & ViT-S/14 rect frozen & LatTF & B & 0.89 & 0.80 & 44.1 \\
F1-M1-d & mae & ViT-S/14 rect trainable & LatTF & B & 0.90 & 0.84 & 18.7 \\
F1-M1-e & mae & ViT-S/14 sq frozen & LatTF & B & 0.84 & 0.70 & 151.7 \\
F1-M1-f & mae & ViT-S/14 sq trainable & LatTF & B & 0.84 & 0.70 & 130.4 \\
F1-M2 & mae & ViT-B/16 frozen & LatTF & LV & 0.92 & 0.79 & 42.0 \\
F1-M2-b & mae & ViT-B/16 trainable & LatTF & LV & 0.89 & 0.77 & 35.2 \\
F1-M2-c & mae & ViT-S/14 rect frozen & LatTF & LV & 0.93 & 0.76 & 49.0 \\
F1-M2-d & mae & ViT-S/14 rect trainable & LatTF & LV & 0.93 & 0.81 & 29.9 \\
F1-M2-e & mae & ViT-S/14 sq frozen & LatTF & LV & 0.89 & 0.65 & 109.6 \\
F1-M2-f & mae & ViT-S/14 sq trainable & LatTF & LV & 0.88 & 0.68 & 98.6 \\
F1-M3 & mae & ViT-B/16 frozen & LatTF & P & 0.89 & 0.87 & 49.6 \\
F1-M3-b & mae & ViT-B/16 trainable & LatTF & P & 0.81 & 0.95 & 4.5 \\
F1-M3-c & mae & ViT-S/14 rect frozen & LatTF & P & 0.89 & 0.88 & 29.3 \\
F1-M3-d & mae & ViT-S/14 rect trainable & LatTF & P & 0.88 & 0.83 & 13.5 \\
F1-M3-e & mae & ViT-S/14 sq frozen & LatTF & P & 0.82 & 0.64 & 162.8 \\
F1-M3-f & mae & ViT-S/14 sq trainable & LatTF & P & 0.87 & 0.72 & 110.7 \\
F1-M4 & mae & ViT-B/16 frozen & LatTF & S & 0.89 & 0.88 & 40.7 \\
F1-M4-b & mae & ViT-B/16 trainable & LatTF & S & 0.77 & 0.95 & 23.1 \\
F1-M4-c & mae & ViT-S/14 rect frozen & LatTF & S & 0.90 & 0.90 & 23.2 \\
F1-M4-d & mae & ViT-S/14 rect trainable & LatTF & S & 0.88 & 0.94 & 6.5 \\
F1-M4-e & mae & ViT-S/14 sq frozen & LatTF & S & 0.82 & 0.70 & 142.8 \\
F1-M4-f & mae & ViT-S/14 sq trainable & LatTF & S & 0.87 & 0.74 & 100.4 \\
\midrule
P1-B1 &  &  &  &  & 0.99 & 0.64 & 93.0 \\
\midrule
P1-B4-b & resnet34 & supervised & TF & S & 0.97 & 0.81 & 67.4 \\
\midrule
P1-D1 & dinov2 & ViT-S/14 frozen & TF & B & 0.58 & 0.83 & 97.5 \\
P1-D1-b & dinov2 & ViT-S/14 trainable & TF & B & 0.54 & 0.91 & 48.8 \\
P1-D1-c & dinov2 & ViT-S/14 rect frozen & TF & B & 0.81 & 0.92 & 20.8 \\
P1-D1-d & dinov2 & ViT-S/14 rect trainable & TF & B & 0.79 & 0.93 & 13.3 \\
P1-D1-e & dinov2 & ViT-S/14 sq frozen & TF & B & 0.75 & 0.77 & 112.7 \\
P1-D1-f & dinov2 & ViT-S/14 sq trainable & TF & B & 0.76 & 0.74 & 103.2 \\
P1-D2 & dinov2 & ViT-S/14 frozen & TF & LV & 0.66 & 0.99 & 3.4 \\
P1-D2-b & dinov2 & ViT-S/14 trainable & TF & LV & 0.43 & 0.95 & 14.6 \\
P1-D2-c & dinov2 & ViT-S/14 rect frozen & TF & LV & 0.82 & 0.90 & 25.2 \\
P1-D2-d & dinov2 & ViT-S/14 rect trainable & TF & LV & 0.80 & 0.93 & 21.2 \\
P1-D2-e & dinov2 & ViT-S/14 sq frozen & TF & LV & 0.80 & 0.75 & 93.4 \\
P1-D2-f & dinov2 & ViT-S/14 sq trainable & TF & LV & 0.80 & 0.71 & 105.2 \\
P1-D3 & dinov2 & ViT-S/14 frozen & TF & P & 0.60 & 0.83 & 98.4 \\
P1-D3-b & dinov2 & ViT-S/14 trainable & TF & P & 0.53 & 0.92 & 61.7 \\
P1-D3-c & dinov2 & ViT-S/14 rect frozen & TF & P & 0.77 & 0.92 & 17.2 \\
P1-D3-d & dinov2 & ViT-S/14 rect trainable & TF & P & 0.77 & 0.92 & 10.3 \\
P1-D3-e & dinov2 & ViT-S/14 sq frozen & TF & P & 0.76 & 0.76 & 121.8 \\
P1-D3-f & dinov2 & ViT-S/14 sq trainable & TF & P & 0.75 & 0.71 & 115.3 \\
P1-D4 & dinov2 & ViT-S/14 frozen & TF & S & 0.60 & 0.79 & 116.6 \\
P1-D4-b & dinov2 & ViT-S/14 trainable & TF & S & 0.56 & 0.88 & 90.4 \\
P1-D4-c & dinov2 & ViT-S/14 rect frozen & TF & S & 0.80 & 0.92 & 18.6 \\
P1-D4-d & dinov2 & ViT-S/14 rect trainable & TF & S & 0.73 & 0.95 & 6.5 \\
P1-D4-e & dinov2 & ViT-S/14 sq frozen & TF & S & 0.75 & 0.78 & 76.8 \\
P1-D4-f & dinov2 & ViT-S/14 sq trainable & TF & S & 0.76 & 0.81 & 62.1 \\
\midrule
P1-I1 & ijepa & ViT-H/14 frozen & TF & B & 0.79 & 0.85 & 29.6 \\
P1-I1-b & ijepa & ViT-H/14 trainable & TF & B & 0.73 & 0.83 & 52.8 \\
P1-I1-c & ijepa & ViT-S/14 rect frozen & TF & B & 0.80 & 0.91 & 15.7 \\
P1-I1-d & ijepa & ViT-S/14 rect trainable & TF & B & 0.78 & 0.91 & 20.1 \\
P1-I1-e & ijepa & ViT-S/14 sq frozen & TF & B & 0.84 & 0.61 & 189.7 \\
P1-I1-f & ijepa & ViT-S/14 sq trainable & TF & B & 0.83 & 0.70 & 143.1 \\
P1-I2 & ijepa & ViT-H/14 frozen & TF & LV & 0.80 & 0.79 & 41.4 \\
P1-I2-b & ijepa & ViT-H/14 trainable & TF & LV & 0.71 & 0.95 & 5.1 \\
P1-I2-c & ijepa & ViT-S/14 rect frozen & TF & LV & 0.76 & 0.92 & 19.3 \\
P1-I2-d & ijepa & ViT-S/14 rect trainable & TF & LV & 0.76 & 0.91 & 21.6 \\
P1-I2-e & ijepa & ViT-S/14 sq frozen & TF & LV & 0.90 & 0.65 & 141.2 \\
P1-I2-f & ijepa & ViT-S/14 sq trainable & TF & LV & 0.85 & 0.72 & 116.8 \\
P1-I3 & ijepa & ViT-H/14 frozen & TF & P & 0.79 & 0.86 & 16.1 \\
P1-I3-b & ijepa & ViT-H/14 trainable & TF & P & 0.71 & 0.87 & 47.6 \\
P1-I3-c & ijepa & ViT-S/14 rect frozen & TF & P & 0.77 & 0.92 & 13.5 \\
P1-I3-d & ijepa & ViT-S/14 rect trainable & TF & P & 0.78 & 0.93 & 13.5 \\
P1-I3-e & ijepa & ViT-S/14 sq frozen & TF & P & 0.84 & 0.62 & 172.3 \\
P1-I3-f & ijepa & ViT-S/14 sq trainable & TF & P & 0.80 & 0.69 & 115.7 \\
P1-I4 & ijepa & ViT-H/14 frozen & TF & S & 0.80 & 0.85 & 12.4 \\
P1-I4-b & ijepa & ViT-H/14 trainable & TF & S & 0.62 & 0.93 & 21.7 \\
P1-I4-c & ijepa & ViT-S/14 rect frozen & TF & S & 0.80 & 0.98 & 5.1 \\
P1-I4-d & ijepa & ViT-S/14 rect trainable & TF & S & 0.81 & 0.98 & 6.0 \\
P1-I4-e & ijepa & ViT-S/14 sq frozen & TF & S & 0.84 & 0.66 & 174.6 \\
P1-I4-f & ijepa & ViT-S/14 sq trainable & TF & S & 0.79 & 0.75 & 115.4 \\
\midrule
P1-M1 & mae & ViT-B/16 frozen & TF & B & 0.88 & 0.86 & 57.1 \\
P1-M1-b & mae & ViT-B/16 trainable & TF & B & 0.66 & 0.96 & 2.2 \\
P1-M1-c & mae & ViT-S/14 rect frozen & TF & B & 0.89 & 0.81 & 38.0 \\
P1-M1-d & mae & ViT-S/14 rect trainable & TF & B & 0.89 & 0.78 & 16.9 \\
P1-M1-e & mae & ViT-S/14 sq frozen & TF & B & 0.83 & 0.64 & 177.3 \\
P1-M1-f & mae & ViT-S/14 sq trainable & TF & B & 0.83 & 0.66 & 144.2 \\
P1-M2 & mae & ViT-B/16 frozen & TF & LV & 0.89 & 0.75 & 61.8 \\
P1-M2-b & mae & ViT-B/16 trainable & TF & LV & 0.79 & 0.88 & 7.8 \\
P1-M2-c & mae & ViT-S/14 rect frozen & TF & LV & 0.93 & 0.78 & 42.8 \\
P1-M2-d & mae & ViT-S/14 rect trainable & TF & LV & 0.93 & 0.83 & 21.6 \\
P1-M2-e & mae & ViT-S/14 sq frozen & TF & LV & 0.87 & 0.63 & 132.6 \\
P1-M2-f & mae & ViT-S/14 sq trainable & TF & LV & 0.85 & 0.63 & 108.1 \\
P1-M3 & mae & ViT-B/16 frozen & TF & P & 0.86 & 0.74 & 66.6 \\
P1-M3-b & mae & ViT-B/16 trainable & TF & P & 0.76 & 0.88 & 35.7 \\
P1-M3-c & mae & ViT-S/14 rect frozen & TF & P & 0.91 & 0.87 & 25.8 \\
P1-M3-d & mae & ViT-S/14 rect trainable & TF & P & 0.85 & 0.80 & 10.9 \\
P1-M3-e & mae & ViT-S/14 sq frozen & TF & P & 0.82 & 0.64 & 165.7 \\
P1-M3-f & mae & ViT-S/14 sq trainable & TF & P & 0.84 & 0.64 & 119.7 \\
P1-M4 & mae & ViT-B/16 frozen & TF & S & 0.88 & 0.80 & 58.3 \\
P1-M4-b & mae & ViT-B/16 trainable & TF & S & 0.80 & 1.00 & 1.2 \\
P1-M4-c & mae & ViT-S/14 rect frozen & TF & S & 0.88 & 0.89 & 21.4 \\
P1-M4-d & mae & ViT-S/14 rect trainable & TF & S & 0.88 & 0.84 & 7.2 \\
P1-M4-e & mae & ViT-S/14 sq frozen & TF & S & 0.82 & 0.65 & 178.2 \\
P1-M4-f & mae & ViT-S/14 sq trainable & TF & S & 0.86 & 0.74 & 113.7 \\
\end{longtable}

